\begin{abstract}
LLM-based assistants that operate software on a user's behalf usually send screenshots and page content to hosted models, retain little from one session to the next, and can declare a task finished when it is not. We present Agentic Pilot, an open-source agent that runs its language and vision models locally through Ollama and drives a Chromium browser through Playwright. The agent is a LangGraph state machine in which observation, planning, execution, verification and recovery are separate nodes, and a task may be reported complete only when a completion predicate evaluated on the live page holds; each evaluation is stored with before/after screenshots and a DOM snapshot. Grounding uses a compact DOM manifest first and a local vision model only when the planner cannot act on the DOM. Calls are routed to small models by role, CAPTCHAs halt the task for the user, and high-risk intents require human approval. We formalize the gate and define a false completion rate. On 35 constructed browser end states, the gate withheld completion in 15 of the 22 states that did not satisfy their goal, all of which a dispatch-only rule would accept, and rejected 1 of 13 satisfied states, at a median cost of 33.8\,ms per decision; the seven misses identify specific missing checks. A code-level audit shows which mechanisms are active on the live path. We also design an extension to files, processes and native applications with accessibility-tree grounding, typed postconditions and deterministic permission tiers, and give an end-to-end evaluation protocol for local hardware.
\end{abstract}

\begin{IEEEkeywords}
autonomous agents, large language models, local inference, privacy-preserving computing, task automation, execution verification, human-in-the-loop
\end{IEEEkeywords}
